% Programming with InitechBase, Chapter 9: The Aging Report (AN-16).  The
% program is Y2KACCT.PRG of the sample diskette, whose body is byte for byte
% the body of harness/diff/dbf_diff/canon/y2k_accounting.prg
% (test-pubs-samir-programs), and its replayed output (examples/p9-aging.dot)
% is that of canon/y2k_accounting.out.  Documented as the shipped program,
% with the date convention of SET CENTURY OFF (CTOD base 1900;
% core/fn_builtins.c fn_ctod_impl; Using InitechBase, SET CENTURY).
\chapter{The Aging Report}
\Gates{test-pubs-examples, test-pubs-samir-programs}
At the end of each month the Collections desk ages the open invoices: for each
invoice it works out how many days have passed since the due date, flags the
invoices more than thirty days past due as overdue, and totals the unpaid
overdue amounts.  \file{Y2KACCT.PRG} is the aging report of the Initech
Accounting System.

\begin{ChapterCovers}
\item The reporting date
\item Date arithmetic
\item The aging report
\item Reading the program
\end{ChapterCovers}

\ATEHead{Preparing for This Chapter}
Start InitechBase with a fresh copy of the sample diskette, and type
\type{USE INVOICE}.

\ATEHead{The Reporting Date}
The report is run as of a reporting date, the memory variable \texttt{ASOF},
which the supervisor sets at the dot prompt before the run.  The report can
therefore be run again for any month-end.  The date is entered, like every date
of the system, in the standard form month, day and year with two digits each,
converted with \fn{CTOD()}.

Under the standard date convention of InitechBase (SET CENTURY OFF), a year
entered with two digits is a year of the nineteen hundreds:
\DotExample{examples/p9-dates}
The due dates of the ledger were entered and stored the same way, on the data
entry screen, so the reporting date and the due dates follow one convention.

\ATEHead{Running the Report}
\begin{Steps}
\item Type \type{ASOF = CTOD('01/31/00')}.
\item Type \type{DO Y2KACCT}.
\end{Steps}
\DotExample{examples/p9-aging}
\SideHead{Age}The age of an invoice is the reporting date minus its due date,
in days.  An invoice whose due date is later than the reporting date has a
negative age: it is not yet due.

\SideHead{Status}An invoice more than thirty days past due is
\texttt{OVERDUE}; every other invoice is \texttt{CURRENT}.

\SideHead{Total}The last line totals the amounts of the overdue invoices that
have not been paid.  For this run every invoice is current, and the total is
zero.

\ATEHead{Reading the Program}
\ProgramListing{Y2KACCT.PRG}
\SideHead{Headings}The first lines print the title and the reporting date,
converted with \fn{DTOC()}.

\SideHead{The loop}For each invoice, \texttt{AGE = ASOF - DUEDATE} subtracts the
dates and gives the age in days.  \fn{IIF()} chooses the status, and \cmd{?}
prints the line, with the age converted by \fn{STR()} to six characters.

\SideHead{The total}\cmd{IF AGE > 30 .AND. .NOT. PAID} adds the amount of each
unpaid overdue invoice to \texttt{OVERTOT}, which is printed after the loop.

\begin{AteNote}
Set \texttt{ASOF} before you run the report.  The report uses the reporting date
throughout, not the date of the computer, so that a month can be closed again
with the same figures.
\end{AteNote}

\ATEHead{Where Do You Go From Here?}
Chapter~10 is the second month-end program: the finance-charge posting.

\ATEHead{Summary}
\begin{enumerate}
\item The reporting date is the memory variable \texttt{ASOF}, set with
  \fn{CTOD()} before the run.
\item Dates entered with two-digit years follow the standard date convention.
\item One date minus another is a number of days; \fn{IIF()} chooses the status.
\end{enumerate}
